% concatenated from main.tex
\documentclass[11pt]{amsart}
\usepackage{amsmath,amscd,amssymb,amsthm,mathtools}
\usepackage[mathscr]{eucal}
\usepackage[frame,cmtip,arrow,matrix,line,graph,curve]{xy}
\usepackage{graphpap, color}
\usepackage{cancel}
\usepackage{verbatim}
\usepackage{adjustbox}
\usepackage[margin=1in]{geometry}
\usepackage{pgfplots}
\pgfplotsset{compat=1.18}

\usepackage{tikz}
\usepackage{float}
\usepackage{booktabs}
\usepackage{multirow}
\usepackage{tabularx}
\usepackage{enumitem}
\usepackage{microtype}


\usepackage{hyperref}
%\usepackage[notcite, notref]{showkeys}

\setlength{\evensidemargin}{\oddsidemargin}



\numberwithin{equation}{section}

\newtheorem{theo}{Theorem}[section]
\newtheorem{coro}[theo]{Corollary}
\newtheorem{prop}[theo]{Proposition}
\newtheorem{lemm}[theo]{Lemma}

\newtheorem{conj}[theo]{Conjecture}
\newtheorem{assu}[theo]{Assumption}

\theoremstyle{definition}
\newtheorem{defi}[theo]{Definition}
\newtheorem{ques}[theo]{Question}
\newtheorem{exam}[theo]{Example}
\newtheorem{rema}[theo]{Remark}
\newtheorem{prob}[theo]{Problem}
\newtheorem{solu}[theo]{Solution}



\def\blue{\textcolor{blue}}
\def\green{\textcolor{green}}
\def\red{\textcolor{red}}
\def\black{\textcolor{black}}
\def\cyan{\textcolor{cyan}}


\newcommand{\A }{\mathbb{A}}
\newcommand{\CC}{\mathbb{C}}
\newcommand{\EE}{\mathbb{E}}
\newcommand{\LL}{\mathbb{L}}
\newcommand{\PP}{\mathbb{P}}
\newcommand{\QQ}{\mathbb{Q}}
\newcommand{\RR}{\mathbb{R}}
\newcommand{\ZZ}{\mathbb{Z}}
\newcommand{\NN}{\mathbb{N}}
\newcommand{\HH}{\mathbb{H}}


%% mathbf

\newcommand{\bL}{\mathbf{L}}
\newcommand{\bT}{\mathbf{T}}
\newcommand{\bk}{\mathbf{k}}
\newcommand{\bm}{\mathbf{m}}
\newcommand{\bn}{\mathbf{n}}
\newcommand{\bp}{\mathbf{p}}
\newcommand{\bq}{\mathbf{q}}
\newcommand{\bs}{\mathbf{s}}
\newcommand{\bt}{\mathbf{t}}
\newcommand{\bw}{\mathbf{w}}
\newcommand{\bx}{\mathbf{x}}
\newcommand{\kk}{\bk}
\def\bA{\mathbf{A}}
\newcommand{\bP}{\mathbf{P}}
\newcommand{\bX}{\mathbf{X}}
\newcommand{\bS}{\mathbf{S}}
\newcommand{\bM}{\mathbf{M}}
\newcommand{\bH}{\mathbf{H}}
\newcommand{\bC}{\mathbf{C}}

%% mathcal

\def\cA{{\mathcal A}}
\def\cB{{\mathcal B}}
\def\cC{{\mathcal C}}
\def\cD{{\mathcal D}}
\def\cE{{\mathcal E}}
\def\cF{{\mathcal F}}
\def\cG{{\mathcal G}}
\def\cH{{\mathcal H}}
\def\cK{{\mathcal K}}
\def\cL{{\mathcal L}}
\def\cM{{\mathcal M}}
\def\cN{{\mathcal N}}
\def\cO{{\mathcal O}}
\def\cP{{\mathcal P}}
\def\cQ{{\mathcal Q}}
\def\cR{{\mathcal R}}
\def\cS{{\mathcal S}}
\def\cT{{\mathcal T}}
\def\cU{{\mathcal U}}
\def\cV{{\mathcal V}}
\def\cW{{\mathcal W}}
\def\cX{{\mathcal X}}
\def\cY{{\mathcal Y}}
\def\cZ{{\mathcal Z}}
\def\cI{{\mathcal I}}
\def\cJ{{\mathcal J}}

%% mathfrak

\def\fB{\mathfrak{B}}
\def\fC{\mathfrak{C}}
\def\fD{\mathfrak{D}}
\def\fE{\mathfrak{E}}
\def\fF{\mathfrak{F}}
\def\fM{\mathfrak{M}}
\def\fV{\mathfrak{V}}
\def\fX{\mathfrak{X}}
\def\fY{\mathfrak{Y}}

\def\ff{\mathfrak{f}}
\def\fm{\mathfrak{m}}
\def\fp{\mathfrak{p}}
\def\ft{\mathfrak{t}}
\def\fu{\mathfrac{u}}
\def\fv{\mathfrak{v}}
\def\frev{\mathfrak{rev}}




%% tilde, English

\newcommand{\tC}{\tilde{C}}
\newcommand{\tD}{\tilde{D}}
\newcommand{\tE}{\tilde{E}}
\newcommand{\tF}{\tilde{F}}
\newcommand{\tK}{\tilde{K}}
\newcommand{\tL}{\tilde{L}}
\newcommand{\tT}{\tilde{T}}
\newcommand{\tY}{\tilde{Y}}



%% hat

\newcommand{\hD}{\hat{D}}
\newcommand{\hL}{\hat{L}}
\newcommand{\hT}{\hat{T}}
\newcommand{\hX}{\hat{X}}
\newcommand{\hY}{\hat{Y}}
\newcommand{\hZ}{\hat{Z}}
\newcommand{\hf}{\hat{f}}
\newcommand{\hu}{\hat{u}}
\newcommand{\hDe}{\hat{Delta}}



\def\begeq{\begin{equation}}
\def\endeq{\end{equation}}
\def\and{\quad{\rm and}\quad}
\def\lra{\longrightarrow }
\def\mapright#1{\,\smash{\mathop{\lra}\limits^{#1}}\,}
\def\mapleft#1{\,\smash{\mathop{\longleftarrow}\limits^{#1}}\,}
\def\sO{{\mathscr O}}
\def\beq{\begin{equation}}
\def\eeq{\end{equation}}
\def\DM{Deligne-Mumford }
\def\wt{\mathrm{wt}}
\def\GG{\mathbb{G} }
\def\Gm{\mathbb{G}_m}
\def\Pic{\mathrm{Pic} }
\def\spec{\mathrm{Spec}\, }
\def\Map{\mathrm{Map} }
\def\oM{\overline{M} }
\def\Hom{\mathrm{Hom} }
\def\bzero{\mathbf{0}}
\def\bone{\mathbf{1}}
\def\bcM{\overline{\cM}}
\def\cAn{\cA^{(n)}}
\def\Yn{Y^{(n)}}
\def\bHn{\bH^{(n)}}


\title{Hodge-Deligne and Poincar\'e polynomials of $\overline{\mathcal M}_{1,n}$} % 1: Betti and Hodge numbers}
\author{Young-Hoon Kiem}
\address{School of Mathematics, Korea Institute for Advanced Study, 85 Hoegiro, Dongdaemun-gu, Seoul 02455, Korea}
\email{kiem@kias.re.kr}

\date{}

\setcounter{tocdepth}{1}
\begin{document}

\begin{abstract}
    By studying wall crossings of Hassett moduli spaces of weighted pointed curves, we prove a recursive formula that reduces the computation of the cohomology of $\bcM_{g,n}$ to that of the moduli space $\bcM_{g,(0^+)^n}$ of stable pointed curves of minimal weights. For $g=1$, we further prove an explicit formula for the Hodge-Deligne and Poincar\'e polynomials of $\bcM_{1,(0^+)^n}$. Combining these, we prove formulas that enable us to compute the cohomology of $\bcM_{1,n}$ efficiently.  
\end{abstract}

\maketitle


\section{Introduction}

The moduli space $\bcM_{g,n}$ of stable curves of genus $g$ with $n$ marked points is one of the most important objects in algebraic geometry \cite{DM, Knu}. It parametrizes connected projective curves $C$ of arithmetic genus $g$ with at worst nodal singularities together with $n$ ordered smooth marked points $p_1,\cdots, p_n$ in $C$ such that 
\begin{enumerate}
\item $p_i$ are all distinct and 
\item $\omega_C(\sum_i p_i)$ is ample. 
\end{enumerate}
It is a smooth proper Deligne-Mumford stack with projective coarse moduli space. It plays a fundamental role in modern enumerative geometry and serves as a model for moduli spaces of higher dimensional varieties. 
For over half a century, much effort has been exerted to understand the geometry and topology of $\bcM_{g,n}$ and there has been a constant stream of new discoveries even until now \cite{BK,BKs,CK,CKL1,CKL2,CKL3,KSY}. 

In this paper, we focus on the cohomology of $\bcM_{1,n}$. 
In principle, the cohomology $H^*(\bcM_{1,n})$ can be computed by \cite{Getzler2} which gives a formula for the $S_n$-equivariant Hodge-Deligne polynomial of $\bcM_{1,n}$. Hence we can compute the ordinary Hodge-Deligne polynomial as well as the $S_n$-action on the cohomology. 
Unfortunately, the formula is rather difficult to use for practical computation: We first need to compute the \emph{equivariant} Hodge-Deligne polynomials of $\cM_{1,n}$, $\cM_{0,n}$ and $\bcM_{0,n}$ for all $n$ and  then expand a power series after applying the plethystic composition. As far as we know, computation for $\bcM_{1,n}$ has been worked out only for $n$ up to 20 \cite{BeT}. 

In contrast, the computation of the Poincar\'e polynomial of $\bcM_{0,n}$ is much easier due to an \emph{explicit recursive} formula \cite[(4.12)]{Manin} which can be proved by Keel's blowup construction \cite{Keel}.
Therefore it seems reasonable to ask the following.
\begin{ques}
    Is there an explicit recursive formula for the Poincar\'e and Hodge-Deligne polynomials of $\bcM_{1,n}$?
\end{ques}
The purpose of this paper is to provide a positive answer to this question so that we get a new and effective way to compute the Hodge and Betti numbers of $\bcM_{1,n}$. 

\medskip

Our computation of $H^*(\bcM_{1,n})$ is achieved in two stages. 
We first reduce the computation to a simpler moduli space $\bcM_{1,(0^+)^n}$, where $0^+$ denotes a positive rational number sufficiently close to $0$.  
We note that many of the strata in $\bcM_{1,n}$ disappear and the computation becomes much easier if we allow all the marked points to coincide. This happens in the moduli space $\bcM_{1,(0^+)^n}$ of stable pointed curves with weight $0^+$ for each marked point \cite{Hst}. 
By studying wall crossings in Hassett's moduli spaces of weighted pointed curves, we prove the following.
\begin{theo}[Theorem \ref{17}] Let $g\ge 1$ and fix a motivic invariant $\bH$ satisfying the blowup formula \eqref{12}. 
Let $\bcM_{g,r|s}=\bcM_{g,(1^r,(0^+)^s)}$ be the moduli space of stable curves of genus $g$ with $r$ points of weight 1 and $s$ points of weight $0^+$. Let $\bH_{r|s}:=\bH_{\bcM_{g,r|s}}$. Then for $r,s\ge 0$ with $r+s>0$ and $n>0$, the following hold:
   \beq\label{32} \bH_{r|s} = \bH_{0|r+s} + \sum_{j=1}^{\min(r,r+s-2)} \sum_{k=0}^{r+s-j-2} \binom{r+s-j}{k}\, \bH_{j|k}  \cdot (\bH_{\PP^{r+s-j-k-1}}-1),\eeq
   \beq\label{33} \bH_{\bcM_{g,n}}= \bH_{n|0} = \bH_{0|n} + \sum_{j=1}^{n-2} \sum_{k=0}^{n-j-2} \binom{n-j}{k}\, \bH_{j|k} \cdot (\bH_{\PP^{n-j-k-1}}-1).\eeq
All $\bH_{r|s}$ are recursively determined by \eqref{32} if we know $\bH_{0|n}=\bH_{\bcM_{g,(0^+)^n}}$ for all $n$.   
\end{theo}

Next we compute the cohomology of $\bcM_{1,(0^+)^n}$ piece by piece. 
Let $$\pi:\bcM_{1,(0^+)^n}\lra \bcM_{1,0^+}=\bcM_{1,1}$$
be the morphism that forgets all the marked points except the first. The inverse image $\cI_n=\pi^{-1}(\cM_{1,1})$ of $\cM_{1,1}$ is the $(n-1)$-fold fiber product of the universal curve $$q:\cC\lra \cM_{1,1}=SL_2(\ZZ)\backslash \HH, \quad \HH=\{z\in \CC\,|\, \mathrm{Im}\, z>0\}$$ and we have
\beq\label{34} R\pi_*\QQ_{\cI_n}=(Rq_*\QQ_\cC)^{\otimes n-1}=\left((\QQ\oplus \QQ[-2])\oplus V[-1]\right)^{\otimes n-1}\eeq
by the relative K\"unneth formula, where $V=R^1q_*\QQ$.
Note that $H^*_c(\cM_{1,1},V^{\otimes k})=0$ for $k$ odd by the action of the generic stabilizer $-I\in SL_2(\ZZ)$.  
By expanding the right hand side of \eqref{34} and applying the Clebsch-Gordan rule, the computation of $H^*_c(\cI_n)$ is reduced to that of $H^*_c(\cM_{1,1}, \mathrm{Sym}^{2j}V)$ which can be obtained by the Eichler-Shimura isomorphism (cf. Lemma \ref{35}). 
Combining these, we find that the Hodge-Deligne polynomial of the interior $\cI_n$ is (cf. Proposition \ref{3})
\beq\label{36}
\sum_{j=0}^{\lfloor\frac{n-1}{2}\rfloor} \binom{n-1}{2j} (uv+1)^{n-1-2j} \left( \alpha_{jj}(uv)^{j+1} - \sum_{i=0}^{j-1} \alpha_{ij}(uv)^i\right),
\eeq
\[-\sum_{j=0}^{\lfloor\frac{n-1}{2}\rfloor} \binom{n-1}{2j} (uv+1)^{n-1-2j}  \sum_{i=0}^{j-1} \alpha_{ij}\beta_{2j-2i+2}(uv)^i(u^{2j-2i+1}+v^{2j-2i+1})  \]
where \[ \alpha_{ij}=\binom{2j}{i}-\binom{2j}{i-1}
\and 
\beta_k=\left\{\begin{matrix}  \lfloor \frac{k}{12}\rfloor -1 & \text{if } k\equiv 2 \text{ mod } 12\\
\lfloor \frac{k}{12}\rfloor & \text{if } k\not\equiv 2 \text{ mod } 12.\end{matrix}\right.
\]

The complement $\cB_n=\pi^{-1}(\infty)=\bcM_{1,n}-\cI_n$ consists of pointed rational nodal curves. 
Let $\cB_{n,j}\subset \bcM_{1,n}$ be the locus of pointed rational curves whose dual graphs are necklaces with $j$ vertices. 
It is obvious that there are $\gamma_{n,j}$
dual graphs with $j$ vertices, where
\[ \gamma_{n,j}=\left\{\begin{matrix} n\\ j\end{matrix}\right\}\frac{(j-1)!}{2} \ \ \text{for }\ \ j>2, \quad \gamma_{n,1}=1, \ \ \gamma_{n,2}=2^{n-1}-1\]
and $\left\{\begin{matrix} n\\ j\end{matrix}\right\}$ denotes the Stirling number of the second kind. 
For $j\ge 2$, $\cB_{n,j}\cong (\CC^*)^{n-j}$ and hence $H^*_c(\cB_{n,j})$ has Hodge-Deligne polynomial $\gamma_{n,j}(uv-1)^{n-j}$. For $j=1,2$, we have to consider an additional action of $S_2$ which acts on $\CC^*$ by $z\mapsto z^{-1}$. Combining all these, we find that the Hodge-Deligne polynomial of the boundary $\cB_n$ is 
\beq\label{37}
\frac{(uv+1)^{n-1}+(uv-1)^{n-1}}{2} + \gamma_{n,2} \frac{(uv+1)^{n-2}+(uv-1)^{n-2}}{2} + \sum_{j=3}^n \gamma_{n,j} (uv-1)^{n-j}.
\eeq
By \eqref{36}, \eqref{37} and  \eqref{32}, we have the following.
\begin{theo} [Theorem \ref{22}]
The Hodge-Deligne polynomial $\bH_{0|n}$ of $\bcM_{1,(0^+)^n}=\bcM_{1,0|n}$ is 
\beq\label{38}
\sum_{j=0}^{\lfloor\frac{n-1}{2}\rfloor} \binom{n-1}{2j} (uv+1)^{n-1-2j} \left( \alpha_{jj}(uv)^{j+1} - \sum_{i=0}^{j-1} \alpha_{ij}(uv)^i\right),
\eeq
\[ +\frac{(uv+1)^{n-1}+(uv-1)^{n-1}}{2} + \gamma_{n,2} \frac{(uv+1)^{n-2}+(uv-1)^{n-2}}{2} + \sum_{j=3}^n \gamma_{n,j} (uv-1)^{n-j} \]
\[-\sum_{j=0}^{\lfloor\frac{n-1}{2}\rfloor} \binom{n-1}{2j} (uv+1)^{n-1-2j}  \sum_{i=0}^{j-1} \alpha_{ij}\beta_{2j-2i+2}(uv)^i(u^{2j-2i+1}+v^{2j-2i+1}).  \]
The Hodge-Deligne polynomial $\bH_{n|0}$ of $\bcM_{1,n}=\bcM_{1,n|0}$ can be computed by the recursive formula \eqref{33} together with \eqref{32} and \eqref{38}.
\end{theo}
Letting $u=v=-t$, we get the Poincar\'e polynomial of $\bcM_{1,n}$. Note that only the last line in \eqref{38} has odd degree terms. 

It is straightforward to implement the above theorem into a computer program and we can compute the Hodge-Deligne and Poincar\'e polynomials of $\bcM_{1,n}$ for large $n$ within a relatively short period of time. %\footnote{With an ordinary laptop, it took less than 5 minutes to compute for $n$ up to 50.} 

\medskip

Motivated by recent discoveries on the cohomology of $\bcM_{0,n}$ \cite{BK, BKs,CK,CKL1,CKL2,CKL3}, we raise the following. 
\begin{ques}
(1)  Can we refine the computations in this paper to prove a new recursive formula for the $S_n$-equivariant Hodge-Deligne polynomial of $\bcM_{1,n}$? 

(2) Is the even degree Poincar\'e polynomial $P^{\mathrm{even}}_{\bcM_{1,n}}(x)=\sum_{i=0}^n x^i\dim H^{2i}(\bcM_{1,n})$ real-rooted?

(3) Do the coefficients of $P^{\mathrm{even}}_{\bcM_{1,n}}(x)$ form a log-concave or $k$-ultra log-concave sequence for some $k$?

(4) Does the probability distribution of the coefficients of $P^{\mathrm{even}}_{\bcM_{1,n}}(x)/P^{\mathrm{even}}_{\bcM_{1,n}}(1)$ converge to a Gaussian distribution? 

(5) Can we find effective estimates of the Betti numbers of $\bcM_{1,n}$ for large $n$? 

(6) Can we also use the method in this paper to calculate the Hodge-Deligne and Poincar\'e polynomials of $\bcM_{g,n}$ for $g\ge 2$? Note that Theorem \ref{17} holds for all $g\ge 1$ and hence it suffices to calculate the Hodge-Deligne polynomial of $\bcM_{g,0|n}$.
\end{ques}



In subsequent papers, we will continue our investigation to seek answers to the above questions. 

\medskip

\noindent\textbf{Convention}. 
All schemes and stacks in this paper are defined over the complex number field $\CC$
and all cohomology groups have rational coefficients. 

\medskip

\noindent\textbf{Acknowledgement}. It is my pleasure to thank Jinwon Choi and Han-Bom Moon for useful discussions. 

\bigskip

\section{Moduli spaces of weighted pointed curves and Hodge-Deligne polynomials}
In this section, we collect necessary facts about Hassett moduli spaces \cite{Hst} and Hodge-Deligne polynomials that will be used throughout this paper. 

\begin{defi} \cite{Hst}
Let $g, n\ge 0$. For $\cA=(a_1,\cdots,a_n)\in (0,1]^n$ with $2g-2+\sum_i a_i>0$, an integral curve $C$ with marked points $p_1, \cdots,p_n\in C$ is called $\cA$-\emph{stable} if \begin{enumerate}
\item $C$ has at worst nodal singularities and $p_i$ are all smooth points;
\item if $p_{i_1}=p_{i_2}=\cdots=p_{i_k}$, then $a_{i_1}+\cdots+a_{i_k}\le 1$;
\item $\omega_C(a_1p_1+\cdots +a_np_n)$ is ample. 
\end{enumerate}
\end{defi}
Let $\bcM_{g,\cA}$ denote the moduli stack of $\cA$-stable $n$-pointed curves of genus $g$. 
In particular, when $a_i>1/2$ for all $i$, $\bcM_{g,\cA}$ coincides with the Knudsen-Deligne-Mumford moduli stack $\bcM_{g,n}$ of stable curves of genus $g$ with $n$ marked points \cite{DM,Knu}.

\medskip

Let us summarize useful facts about $\bcM_{g,\cA}$ from \cite{Hst}.
\begin{theo}\label{2} 
Let $\cA=(a_1,\cdots,a_n)\in (0,1]^n$ with $2g-2+\sum_{i=1}^n a_i>0$.

(1) $\bcM_{g,\cA}$ is a smooth connected proper Deligne-Mumford stack of dimension $3g-3+n$ with projective coarse moduli space.  

(2) If $\cA=(a_1,\cdots,a_n)$, $\cA'=(a'_1,\cdots, a'_n)$ and $a_i\ge a'_i$ for all $i$, then we have a natural morphism 
\beq\label{1} \rho_{\cA,\cA'}:\bcM_{g,\cA}\lra \bcM_{g,\cA'}\eeq 
which is obtained by contracting rational components in an $\cA$-stable curve which are not $\cA'$-stable. 

(3) For $2<r\le n$, if $a'_{i_1}+\cdots+a'_{i_r}\le 1$  and $a_{i_1}+\cdots+a_{i_r}>1$ but any proper subset  of $\{a_{i_1},\cdots, a_{i_r}\}$ has sum at most one while $a_j=a_j'$ for $j\ne i_1,\cdots,i_r$, then the morphism $\rho_{\cA,\cA'}$ is the smooth blowup along the locus $\bcM_{g,\bar \cA'}$ in $\bcM_{g,\cA'}$ where $p_{i_1}=\cdots=p_{i_r}$ and $\bar \cA'$ is obtained from $\cA'$ by replacing $a'_{i_1}, \cdots, a'_{i_r}$ by their sum. 

(4) If $r\le 2$ for any $I\subset [n]$ such that $a_{i_1}+\cdots+a_{i_r}>1$ and $a'_{i_1}+\cdots+a'_{i_r}\le 1$, then $\rho_{\cA,\cA'}$ is an isomorphism. 
\end{theo}
\begin{proof}
    (1) is \cite[Theorem 2.1]{Hst} and (2) is \cite[Theorem 4.1]{Hst} while (3) is \cite[Remark 4.6]{Hst} and (4) is \cite[Corollary 4.7]{Hst}. 
\end{proof}

\medskip

Let $\bH$ be a motivic invariant with values in an integral domain $R$ which assigns 
to each variety $X$ over $\CC$ an element $\bH_X\in R$ such that  
\begin{enumerate}
    \item $\bH_X=\bH_{X-Z}+\bH_Z$ for any closed $Z\subset X$ and
    \item $\bH_{X\times Y}=\bH_X\cdot \bH_Y.$
\end{enumerate}
If $X\to Y$ is a Zariski locally trivial morphism with $Y$ connected, then $\bH_X=\bH_Y\cdot \bH_F$ where $F$ denotes a fiber. 

In this paper, we will only consider motivic invariants satisfying the following \emph{blowup formula}: If $X\to Y$ is the smooth blowup along $Z$ of codimension $c$, then %we have the blowup formula
\beq\label{12}
\bH_X=\bH_Y + \bH_Z \cdot (\bH_{\PP^{c-1}}-1).\eeq 

Our main example is the Hodge-Deligne polynomial
$$\bH_X(u,v)=\sum_k (-1)^k\sum_{p,q}u^pv^q\dim H^{p,q;k}_c(X) \in \ZZ[u,v], \quad H^{p,q;k}_c(X)=\mathrm{gr}^p_F\mathrm{gr}_{p+q}^W H^k_c(X)$$
where $\mathrm{gr}_F^p$ and $\mathrm{gr}^W_{p+q}$ denote the graded parts with respect to the Hodge filtration and weight filtration respectively. 
By letting $u=v=-t$, we obtain the virtual Poincar\'e polynomial
$$P_X(t)=\bH_X(-t,-t)\in \ZZ[t]$$
of $X$ which is also a motivic invariant.

%In the following, we will mostly compute the (virtual) Poincar\'e polynomials but the same computation will produce the Hodge(-Deligne) polynomials. 

\bigskip

\section{Reduction to minimally weighted curves}
In this section, we prove an explicit formula (Theorem \ref{17}) that calculates the motivic invariants of $\bcM_{g,n}=\bcM_{g,1^n}$ from those of $\bcM_{g,(0^+)^n}$ where $(0^+)^n\in (0,1]^n$ denotes an element sufficiently close to zero, for $g\ge 1$. 

Fix $g\ge 1$ and $n\ge 1$. Then $\bcM_{g,\cA}$ is a nonempty proper smooth Deligne-Mumford stack for any $\cA\in (0,1]^n$. 
For $0\le r< n$ and $0\le s\le n-r$, let
\beq\label{8} \cAn_{r,s}=(1,\cdots,1,\frac{s}{n-r},\frac1{n-r},\cdots,\frac1{n-r})\in (0,1]^n\eeq
where $1$ is repeated $r-1$ times and $\frac1{n-r}$ is repeated $n-r$ times. 
For instance, $\cAn_{0,s}=(\frac1n,\cdots,\frac1n)$ for any $s$. 
Let 
\beq\label{9} \Yn_{r,s}:=\bcM_{g,\cAn_{r,s}}.\eeq 
In particular, $\Yn_{0,n}=\bcM_{g,(0^+)^n}$ and 
\beq\label{11} \Yn_{n-2,2}=\Yn_{n-1,1}=\bcM_{g,n}.\eeq
Moreover, by Theorem \ref{2} (4),  
\beq\label{15} \Yn_{r,n-r-1}=\Yn_{r,n-r}=\bcM_{g,(1^r,(0^+)^{n-r})}=:\bcM_{g,r|n-r}\eeq
%\ \ \text{ for }\ \ 1\le r<n$$
is the moduli space of stable curves of genus $g$ with $r$ heavy points and $n-r$ light points \cite{KSY}. 

Note that if $r>r'$ or $r=r'$ and $s>s'$, then $\cAn_{r,s}- \cAn_{r',s'}\ge 0$ and hence we have the natural morphism
$$\rho_{(r,s),(r',s')}:\Yn_{r,s}\lra \Yn_{r',s'}$$
by Theorem \ref{2} (2). 

Fix a motivic invariant $\bH$ and let 
\beq\label{14} \bH_{r|s}:=\bH_{\bcM_{g,r|s}}=\bH_{Y^{(r+s)}_{r,s}}\eeq 
be the motivic invariant %Hodge-Deligne polynomial 
of $\bcM_{g,r|s}$. 
Then we have $\bH_{0|n}=\bH_{\bcM_{g,(0^+)^n}}$ and $\bH_{n|0}=\bH_{\bcM_{g,n}}$ by \eqref{15}.

Now our goal is to reach $\Yn_{n-2,2}=\bcM_{g,n}$ starting from $\Yn_{0,n}=\bcM_{g,(0^+)^n}$ by smooth blowups so that we can reduce the computation of $\bH_{n|0}$ to $\bH_{0|n}$. Let us first reach $\Yn_{1,n-1}$. 
The natural map
$$\rho_{\cAn_{1,1},\cAn_{0,n}}:\Yn_{1,1}\lra \Yn_{0,n}$$
is the smooth blowup along the locus $\bcM_{g,1|0}=\bcM_{g,1}$ of $p_1=p_2=\cdots =p_n$ by Theorem \ref{2} (3). 
Likewise, for $1\le k\le n-3$, the natural map 
$$\rho_{\cAn_{1,k+1},\cAn_{1,k}}:\Yn_{1,k+1}\lra \Yn_{1,k}$$
is the smooth blowup along the locus of $p_1=p_{i_1}=\cdots=p_{i_{n-k-1}}$ for $\{i_1,\cdots,i_{n-k-1}\}\subset \{2,\cdots,n\}$ which consists of $\binom{n-1}{k}$ copies of $\bcM_{g,1|k}$.
Note that $\Yn_{1,n-2}=\Yn_{1,n-1}$ by \eqref{15}. 
By applying the blowup formula \eqref{12} to these morphisms, we find that 
\beq\label{13}
\bH_{1|n-1}=\bH_{0|n}+\sum_{k=0}^{n-3} \binom{n-1}{k} \bH_{1|k} (\bH_{\PP^{n-k-2}}-1).
\eeq

Next, let us reach $\Yn_{2,n-2}$ from $\Yn_{1,n-1}$. 
The natural map
$$\rho_{\cAn_{2,1},\cAn_{1,n-1}}:\Yn_{2,1}\lra \Yn_{1,n-1}$$
is the smooth blowup along the locus $\bcM_{g,2|0}=\bcM_{g,2}$ of $p_2=p_3=\cdots =p_n$ by Theorem \ref{2} (3). 
For $1\le k\le n-4$, the natural map 
$$\rho_{\cAn_{2,k+1},\cAn_{2,k}}:\Yn_{2,k+1}\lra \Yn_{2,k}$$
is the smooth blowup along the locus of $p_2=p_{i_1}=\cdots=p_{i_{n-k-2}}$ for $\{i_1,\cdots,i_{n-k-2}\}\subset \{3,\cdots,n\}$ which consists of $\binom{n-2}{k}$ copies of $\bcM_{g,2|k}$.
We have $\Yn_{2,n-3}=\Yn_{1,n-2}$ by \eqref{15} and the blowup formula \eqref{12} gives us 
\beq\label{16}
\bH_{2|n-2}=\bH_{1|n-1}+\sum_{k=0}^{n-4} \binom{n-2}{k} \bH_{2|k} (\bH_{\PP^{n-k-3}}-1).
\eeq
Repeating this way, we obtain the following. 
\begin{theo}\label{17} For $r,s\ge 0$ with $r+s>0$ and $n>0$, the following hold:
   \beq\label{18} \bH_{r|s} = \bH_{0|r+s} + \sum_{j=1}^{\min(r,r+s-2)} \sum_{k=0}^{r+s-j-2} \binom{r+s-j}{k}\, \bH_{j|k}  \cdot (\bH_{\PP^{r+s-j-k-1}}-1).\eeq
   \beq\label{19} \bH_{\bcM_{g,n}}= \bH_{n|0} = \bH_{0|n} + \sum_{j=1}^{n-2} \sum_{k=0}^{n-j-2} \binom{n-j}{k}\, \bH_{j|k} \cdot (\bH_{\PP^{n-j-k-1}}-1).\eeq
All $\bH_{r|s}$ are recursively determined by \eqref{18} if we know $\bH_{0|n}=\bH_{\bcM_{g,(0^+)^n}}$ for all $n$.   
\end{theo}
In particular, the motivic invariants $\bH_{\bcM_{g,n}}=\bH_{n|0}$ can be recursively calculated by \eqref{18} if we know the motivic invariants $\bH_{\bcM_{g,(0^+)^n}}=\bH_{0|n}$ for all $n\ge 1$. 
Conversely, if we know $\bH_{\bcM_{g,n}}=\bH_{n|0}$ for all $n\ge 1$, we can recursively calculate $\bH_{\bcM_{g,(0^+)^n}}=\bH_{0|n}$ for all positive $n$ by \eqref{18}.





\bigskip

\section{Moduli spaces of minimally weighted curves}
In this section, we let $g=1$ and let $\bH$ be the Hodge-Deligne polynomial. 
We find an explicit formula that calculates the Hodge-Deligne polynomials of the moduli spaces $\bcM_{1,(0^+)^n}$
of stable curves of genus 1 with $n$ marked points whose weights are $\le  1/n$. By Theorem \ref{17}, we can then calculate $\bH_{\bcM_{1,n}}$
for all $n>0$.

Note that in $\bcM_{1,(0^+)^n}$, all the marked points are allowed to coincide as long as they stay away from the nodes. 
Let 
$$\pi=\pi_n:\bcM_{1,(0^+)^n}\lra \bcM_{1,0^+}=\bcM_{1,1}$$
be the morphism that forgets all the marked points except the first. Using the stratification 
$$\bcM_{1,1}=\cM_{1,1}\sqcup \{\infty\}$$
where $\infty$ parameterizes the unique singular rational curve, we have the decomposition
$$\bcM_{1,(0^+)^n}=\cI_n\sqcup \cB_n, \quad \cB_n=\pi_n^{-1}(\infty), \ \ \cI_n=\pi_n^{-1}(\cM_{1,1})$$
into the interior and the boundary. 



Let $I_n=\bH_{\cI_n}$ and $B_n=\bH_{\cB_n}$ be the Hodge-Deligne polynomials of $\cI_n$ and $\cB_n$ respectively so that 
$$\bH_{0|n}= \bH_{\bcM_{1,(0^+)^n}}=I_n+B_n $$
is the Hodge-Deligne polynomial of $\bcM_{1,(0^+)^n}$.

\subsection{The Hodge-Deligne Polynomial $I_n$ of $\cI_n$} 
In this subsection, we prove the following. 
\begin{prop}\label{3} $I_n(u,v)=I_n^{\mathrm{even}}(u,v)+I^{\mathrm{odd}}_n(u,v)$ with  
\[
I_n^{\mathrm{even}}(u,v) = 
\sum_{j=0}^{\lfloor\frac{n-1}{2}\rfloor} \binom{n-1}{2j} (uv+1)^{n-1-2j} \left( \alpha_{jj}(uv)^{j+1} - \sum_{i=0}^{j-1} \alpha_{ij}(uv)^i\right),
\]
\[I_n^{\mathrm{odd}}(u,v)=-\sum_{j=0}^{\lfloor\frac{n-1}{2}\rfloor} \binom{n-1}{2j} (uv+1)^{n-1-2j}  \sum_{i=0}^{j-1} \alpha_{ij}\beta_{2j-2i+2}(uv)^i(u^{2j-2i+1}+v^{2j-2i+1})  \]
where \beq\label{21} \alpha_{ij}=\binom{2j}{i}-\binom{2j}{i-1}
\and 
\beta_k=\left\{\begin{matrix}  \lfloor \frac{k}{12}\rfloor -1 & \text{if } k\equiv 2 \text{ mod } 12\\
\lfloor \frac{k}{12}\rfloor & \text{if } k\not\equiv 2 \text{ mod } 12 .\end{matrix}\right.
\eeq
\end{prop}
See Example \ref{5} for an explicit computation of $I_n$ for $n\le 10$. 

\medskip

Let $q:\cC\to \cM_{1,1}$ denote the universal curve.
By our choice of stability condition, all the marked points are allowed to coincide and hence $\cI_n$ is nothing but the $(n-1)$-fold fiber product 
$$\pi:\cI_n=\cC\times_{\cM_{1,1}}\cC\times_{\cM_{1,1}}\cdots\times_{\cM_{1,1}}\cC\lra \cM_{1,1}$$
which is a smooth projective fibration with fiber $E^{n-1}$ for an elliptic curve $E\in \cM_{1,1}$.  

Let $V=R^1q_*\QQ$. Then $V$ is a local system of rank 2 over $\cM_{1,1}$ and $$Rq_*\QQ\cong (\QQ\oplus\QQ[-2])\oplus V[-1].$$ 
Let $\pi'=\pi|_{\cI_n}$. By the relative K\"unneth theorem, we have 
\beq\label{25}
R\pi'_*\QQ_{\cI_n}\cong ((\QQ\oplus \QQ[-2])\oplus V[-1])^{\otimes n-1}
=\sum_{k=0}^{n-1} \binom{n-1}{k} C_k\otimes V^{\otimes k}[-k]
\eeq
where $C_k=(\QQ\oplus \QQ[-2])^{\otimes n-1-k}$.

Since $\cM_{1,1}\cong SL_2(\ZZ)\backslash \HH$ where $\HH$ denotes the upper half plane and the generic stabilizer $-I\in SL_2(\ZZ)$ acts on $V^{\otimes k}$ as $(-1)^k$, we have the vanishing 
$$H^*_c(\cM_{1,1},V^{\otimes k})=0, \quad \forall \text{odd } k.$$  
Hence \eqref{25} gives us
\beq\label{26}
H^*_c(\cI_n)=H^*_c(\cM_{1,1},R\pi'_*\QQ_{\cI_n})
\cong \sum_{j=0}^{\lfloor \frac{n-1}{2}\rfloor} \binom{n-1}{2j}  C_{2j}\otimes H^*_c(\cM_{1,1},V^{\otimes 2j})[-2j].\eeq

By the Clebsch-Gordan rule and \eqref{21}, we have 
$$V^{\otimes 2j}\cong \sum_{i=0}^{j-1} \alpha_{ij} V_{2j-2i}[2j-2i]\oplus \alpha_{jj}\QQ, \quad \text{where} \ \ V_{2k}=\mathrm{Sym}^{2k}V.$$
Therefore \eqref{26} gives us 
\beq\label{27}
H^*_c(\cI_n)\cong 
\sum_{j=0}^{\lfloor \frac{n-1}{2}\rfloor} \binom{n-1}{2j}  C_{2j}\otimes 
\left( \alpha_{jj}H^*_c(\cM_{1,1},\QQ)[-2j]+\sum_{i=0}^{j-1} \alpha_{ij} H^*_c(\cM_{1,1},V_{2j-2i})[-2i]\right).\eeq
Since $H^*_c(\cM_{1,1},\QQ)\cong \QQ[-2]$, Proposition \ref{3} follows immediately from the following.

\begin{lemm}\label{35} For $k>0$, we have an isomorphism
    $$H^*_c(\cM_{1,1},V_{2k})=H^1_c(\cM_{1,1},V_{2k})\cong \QQ\oplus (S_{2k+2}\oplus \bar S_{2k+2})$$
where $S_{2k+2}$ denotes the space of cusp forms of weight $2k+2$ with respect to $SL_2(\ZZ)$. The first factor $\QQ$ has weight $0$ and  $S_{2k+2}$ is of Hodge type $(2k+1,0)$. 
\end{lemm}
\begin{proof}
    Let $\jmath:\cM_{1,1}\to \bcM_{1,1}$ denote the inclusion map. Then we have an exact triangle
    $$\jmath_!V_{2k}\lra \jmath_*V_{2k}\lra \jmath_*V_{2k}|_\infty\lra .$$
    By the Eichler-Shimura isomorphism \cite[Theorem 8.4]{Shi}, we have 
    $$H^*(\bcM_{1,1},\jmath_*V_{2k})=H^1(\bcM_{1,1},\jmath_*V_{2k})\cong H^1_{\mathrm{par}}(SL_2(\ZZ),\mathrm{Sym}^{2k}(\CC^2))\cong S_{2k+2}\oplus \bar S_{2k+2}$$ where $S_{2k+2}$ has Hodge type $(2k+1,0)$. 
    We thus find that $H^*_c(\cM_{1,1},V_{2k})=H^1_c(\cM_{1,1},V_{2k})$ and have an exact sequence
\[    0\lra H^0(\infty,\jmath_*V_{2k}|_\infty)\lra H^1_c(\cM_{1,1},V_{2k})\lra S_{2k+2}\oplus \bar S_{2k+2} \lra 0 .\]
As the cusp $\infty$ has the stabilizer group generated by
\[\sigma:=\left( \begin{matrix}
1&1\\ 0&1    
\end{matrix}\right),\]
$H^0(\infty,\jmath_*V_{2k}|_\infty)$ is the $\sigma$-invariant part $(\mathrm{Sym}^{2k}\QQ^2)^\sigma\cong \QQ$. This completes the proof. 
\end{proof}
It is well known that $\dim S_{2k+2}=\beta_{2k+2}$ for $k>0$ by Riemann-Roch, using the notation of \eqref{21}.  

\medskip





\subsection{The Hodge-Deligne Polynomial $B_n$ of $\cB_n$}

Let $\gamma_{n,1}=1$, $\gamma_{n,2}=2^{n-1}-1$ and for $k>2$, let \beq\label{20} \gamma_{n,k} =\left\{\begin{matrix} n\\ k\end{matrix}\right\}\frac{(k-1)!}{2} = \frac{1}{2k} \sum_{i=0}^k (-1)^{k-i} \binom{k}{i} i^n\eeq
where $\left\{\begin{matrix} n\\ k\end{matrix}\right\}$ denotes the Stirling number of the second kind, i.e. the number of ways to partition $[n]=\{1,2,\cdots,n\}$ into $k$ unordered subsets. 
%Let $$f_m(x) = \frac{(x+1)^m+(x-1)^m}{2} \and g_m(x) = (x-1)^m.$$ 
In this subsection, we prove the following. 
\begin{prop}\label{4} %$B_n(u,v)$ is uniquely determined by 
With $\gamma_{n,k}$ in \eqref{20}, 
\[
B_n(u,v) = \frac{(uv+1)^{n-1}+(uv-1)^{n-1}}{2} + \gamma_{n,2} \frac{(uv+1)^{n-2}+(uv-1)^{n-2}}{2} + \sum_{j=3}^n \gamma_{n,j} (uv-1)^{n-j}
\]
\end{prop}
See Example \ref{6} for the calculation of $B_n$ for $n\le 10$. 

\medskip

Observe that we have a decomposition $$\cB_n=\bigsqcup_{1\le j\le n}\cB_{n,j}$$ where 
$\cB_{n,j}$ denotes the locus of curves with $j$ rational components i.e. the dual graph is a necklace with $j$ vertices.

If $(C,p_1,\cdots, p_n)\in \cB_{n,1}$, $C$ is a rational curve with  normalization $\PP^1$. Let $0,\infty\in \PP^1$ map to the node of $C$. Then 
we find that 
$$\cB_{n,1}\cong (\CC^*)^n/(\CC^*\rtimes S_2)\cong (\CC^*)^{n-1}/S_2$$
where $-1\in S_2$ acts by $(z_1,\cdots,z_{n-1})\mapsto (z_1^{-1},\cdots,z_{n-1}^{-1}).$
Therefore, we have $$H^*_c(\cB_{n,1})\cong (H^*_c(\CC^*)^{\otimes n-1})^{S_2}$$
whose Hodge-Deligne polynomial is 
\beq\label{28}
\bH_{\cB_{n,1}}(u,v)=\frac12\left( (uv+1)^{n-1} + (uv-1)^{n-1} \right).\eeq

If $(C,p_1,\cdots,p_n)\in \cB_{n,2}$, $C=C_1\cup C_2$ with $C_1,C_2\cong \PP^1$.
The locus $\cB_{n,2,k}$ of such pointed curves where one component has $k$ marked points while the other has $n-k$ points is 
$$((\CC^*)^k/\CC^*\times (\CC^*)^{n-k}/\CC^*)/S_2\cong (\CC^*)^{n-2}/S_2$$
where $S_2$ interchanges the two nodes. We thus find that 
\beq\label{29}
\bH_{\cB_{n,2}}(u,v)=(2^{n-1}-1) \frac12\left( (uv+1)^{n-2} + (uv-1)^{n-2} \right).\eeq

For $j\ge 3$, the number of ways to distribute $n$ marked points to $j$ components is $\gamma_{n,j}$ in \eqref{20}. For each such distribution, we have a connected component of $\cB_{n,j}$ isomorphic to $(\CC^*)^{n-j}$. Hence we find that 
\beq\label{30}
\bH_{\cB_{n,j}}(u,v) = \gamma_{n,j}(uv-1)^{n-j},\quad \text{for } j\ge 3.\eeq

Adding \eqref{28}, \eqref{29} and \eqref{30}, we obtain Proposition \ref{4}. 
 
\medskip

\subsection{The Hodge-Deligne polynomial $\bH_{\bcM_{1,(0^+)^n}}$} 

By Propositions \ref{3} and \ref{4}, we obtain the following.
\begin{theo}\label{22}
The Hodge-Deligne polynomial $\bH_{\bcM_{1,(0^+)^n}}=\bH_{0|n}$ of $\bcM_{1,(0^+)^n}$ is 
\[ \bH_{0|n}(u,v)= \sum_{j=0}^{\lfloor\frac{n-1}{2}\rfloor} \binom{n-1}{2j} (uv+1)^{n-1-2j} \left( \alpha_{jj}(uv)^{j+1} - \sum_{i=0}^{j-1} \alpha_{ij}(uv)^i\right)
\]
\[ +\frac{(uv+1)^{n-1}+(uv-1)^{n-1}}{2} + \gamma_{n,2} \frac{(uv+1)^{n-2}+(uv-1)^{n-2}}{2} + \sum_{j=3}^n \gamma_{n,j} (uv-1)^{n-j} \]
\[ -\sum_{j=0}^{\lfloor\frac{n-1}{2}\rfloor} \binom{n-1}{2j} (uv+1)^{n-1-2j}  \sum_{i=0}^{j-1} \alpha_{ij}\beta_{2j-2i+2}(uv)^i(u^{2j-2i+1}+v^{2j-2i+1}) \]
where $\alpha_{ij}$, $\beta_k$ and $\gamma_{n,k}$ are constants defined by \eqref{21} and \eqref{20} respectively. 
\end{theo}

%By Theorem \ref{2} (1), $T_n(x)$ should be palindromic.  
%Combining Examples \ref{5} and \ref{6}, we obtain the following palindromic polynomials. 

See Example \ref{7} for a computation of $\bH_{0|n}$ for $n\le 10$.

By Theorems \ref{22} and \ref{17}, we can now calculate $\bH_{\bcM_{1,n}}(u,v)$ recursively for all $n$.


%\section{Wall crossings to higher weights}
%Letting $P_{0,n}(x) = T_n(x)$, we construct $P_{r,s}(x)$ via the recurrence:
%\[P_{r,s}(x) = P_{0,r+s}(x) + \sum_{j=1}^{\min(r,r+s-2)} \sum_{k=0}^{r+s-j-2} P_{j,k}(x) \binom{r+s-j}{k} \sum_{i=1}^{r+s-j-k-1} x^i\]







%\[P_{n,0}(x) = P_{0,n}(x) + \sum_{j=1}^{n-2} \sum_{k=0}^{n-j-2} P_{j,k}(x) \binom{n-j}{k} \sum_{i=1}^{n-j-k-1} x^i\]








\bigskip

\section{Computations for $n\le 10$}
%For $k< 12$, $\beta_k=0$ and hence $I_n^{\mathrm{odd}}=0$ for $n\le 10$.  
Letting $x=uv$, by Proposition \ref{3}, we can evaluate $I_n$ as follows. 
\begin{exam}\label{5}
\begin{align*}
I_1(x) &= x \\
I_2(x) &= x^2 + x \\
I_3(x) &= x^3 + 3x^2 + x - 1 \\
I_4(x) &= x^4 + 6x^3 + 6x^2 - 2x - 3 \\
I_5(x) &= x^5 + 10x^4 + 20x^3 + 4x^2 - 14x - 7 \\
I_6(x) &= x^6 + 15x^5 + 50x^4 + 40x^3 - 30x^2 - 49x - 15 \\
I_7(x) &= x^7 + 21x^6 + 105x^5 + 160x^4 - 183x^2 - 139x - 31 \\
I_8(x) &= x^8 + 28x^7 + 196x^6 + 469x^5 + 280x^4 - 427x^3 - 700x^2 - 356x - 63 \\
I_9(x) &= x^9 + 36x^8 + 336x^7 + 1148x^6 + 1386x^5 - 406x^4 - 2436x^3 - 2224x^2 - 860x - 127 \\
I_{10}(x) &= x^{10} + 45x^9 + 540x^8 + 2484x^7 + 4662x^6 + 1764x^5 - 6090x^4 - 9804x^3 - 6372x^2 - 2003x - 255
\end{align*}
\end{exam}

\medskip

By Proposition \ref{4}, we can evaluate $B_n$ as follows with $x=uv$. 
\begin{exam}\label{6}
\begin{align*}
B_1(x) &= 1 \\
B_2(x) &= x + 1 \\
B_3(x) &= x^2 + 3x + 2 \\
B_4(x) &= x^3 + 7x^2 + 9x + 4 \\
B_5(x) &= x^4 + 15x^3 + 31x^2 + 25x + 8 \\
B_6(x) &= x^5 + 31x^4 + 100x^3 + 111x^2 + 65x + 16 \\
B_7(x) &= x^6 + 63x^5 + 316x^4 + 476x^3 + 351x^2 + 161x + 32 \\
B_8(x) &= x^7 + 127x^6 + 987x^5 + 2178x^4 + 1883x^3 + 1023x^2 + 385x + 64 \\
B_9(x) &= x^8 + 255x^7 + 3053x^6 + 10515x^5 + 12307x^4 + 6637x^3 + 2815x^2 + 897x + 128 \\
B_{10}(x) &= x^9 + 511x^8 + 9366x^7 + 51313x^6 + 92466x^5 + 62065x^4 + 21654x^3 + 7423x^2 + 2049x + 256
\end{align*}
\end{exam}

\medskip 

By adding $I_n$ and $B_n$, we find $T_n=\bH_{\bcM_{1,0|n}}$ as follows. 
\begin{exam}\label{7}
\begin{align*}
T_1(x) &= x + 1 \\
T_2(x) &= x^2 + 2x + 1 \\
T_3(x) &= x^3 + 4x^2 + 4x + 1 \\
T_4(x) &= x^4 + 7x^3 + 13x^2 + 7x + 1 \\
T_5(x) &= x^5 + 11x^4 + 35x^3 + 35x^2 + 11x + 1 \\
T_6(x) &= x^6 + 16x^5 + 81x^4 + 140x^3 + 81x^2 + 16x + 1 \\
T_7(x) &= x^7 + 22x^6 + 168x^5 + 476x^4 + 476x^3 + 168x^2 + 22x + 1 \\
T_8(x) &= x^8 + 29x^7 + 323x^6 + 1456x^5 + 2458x^4 + 1456x^3 + 323x^2 + 29x + 1 \\
T_9(x) &= x^9 + 37x^8 + 591x^7 + 4201x^6 + 11901x^5 + 11901x^4 + 4201x^3 + 591x^2 + 37x + 1 \\
T_{10}(x) &= x^{10} + 46x^9 + 1051x^8 + 11850x^7 + 55975x^6 + 94230x^5 + 55975x^4 + 11850x^3 + 1051x^2 + 46x + 1
\end{align*}
\end{exam}
Note that this matches perfectly with Table 2 in \cite{KSY} with $x=uv$.

\medskip

By \eqref{18} and Theorem \ref{22}, we can compute $\bH_{r|s}=\bH_{\bcM_{1,r|s}}$ as follows with $x=uv$. 
\begin{exam}\label{23}
%\noindent \textbf{For $r+s = 1$:}
\begin{align*}
\bH_{0|1} = \bH_{1|0} &= x + 1 \\
%\end{align*}
%\noindent \textbf{For $r+s = 2$:}
%\begin{align*}
\bH_{0|2} = \bH_{1|1} = \bH_{2|0} &= x^2 + 2x + 1 \\
%\end{align*}
%\noindent \textbf{For $r+s = 3$:}
%\begin{align*}
\bH_{0|3} &= x^3 + 4x^2 + 4x + 1 \\
\bH_{1|2} = \bH_{2|1} = \bH_{3|0} &= x^3 + 5x^2 + 5x + 1 \\
%\end{align*}
%\noindent \textbf{For $r+s = 4$:}
%\begin{align*}
\bH_{0|4} &= x^4 + 7x^3 + 13x^2 + 7x + 1 \\
\bH_{1,3} &= x^4 + 11x^3 + 21x^2 + 11x + 1 \\
\bH_{2|2} = \bH_{3|1} = \bH_{4|0} &= x^4 + 12x^3 + 23x^2 + 12x + 1 \\
%\end{align*}
%\noindent \textbf{For $r+s = 5$:}
%\begin{align*}
\bH_{0|5} &= x^5 + 11x^4 + 35x^3 + 35x^2 + 11x + 1 \\
\bH_{1|4} &= x^5 + 22x^4 + 79x^3 + 79x^2 + 22x + 1 \\
\bH_{2|3} &= x^5 + 26x^4 + 97x^3 + 97x^2 + 26x + 1 \\
\bH_{3|2} = \bH_{4|1} = \bH_{5|0} &= x^5 + 27x^4 + 102x^3 + 102x^2 + 27x + 1.
\end{align*}
\end{exam}

Note that this matches perfectly with Table 1 in \cite{KSY} after applying the hook-length formula to extract the Hodge-Deligne polynomials.  

\medskip

By \eqref{19} and Theorem \ref{22}, we can compute $\bH_{n|0}=\bH_{\bcM_{1,n}}$ as follows with $x=uv$. 
\begin{exam}\label{24}
\begin{align*}
\bH_{\bcM_{1,1}} &= x + 1 \\
\bH_{\bcM_{1,2}} &= x^2 + 2x + 1 \\
\bH_{\bcM_{1,3}} &= x^3 + 5x^2 + 5x + 1 \\
\bH_{\bcM_{1,4}} &= x^4 + 12x^3 + 23x^2 + 12x + 1 \\
\bH_{\bcM_{1,5}} &= x^5 + 27x^4 + 102x^3 + 102x^2 + 27x + 1 \\
\bH_{\bcM_{1,6}} &= x^6 + 58x^5 + 421x^4 + 756x^3 + 421x^2 + 58x + 1 \\
\bH_{\bcM_{1,7}} &= x^7 + 121x^6 + 1612x^5 + 5077x^4 + 5077x^3 + 1612x^2 + 121x + 1 \\
\bH_{\bcM_{1,8}} &= x^8 + 248x^7 + 5802x^6 + 31072x^5 + 52402x^4 + 31072x^3 + 5802x^2 + 248x + 1 \\
\bH_{\bcM_{1,9}} &= x^9 + 503x^8 + 19925x^7 + 175036x^6 + 480097x^5 + 480097x^4 + 175036x^3 + 19925x^2 + 503x + 1 \\
\bH_{\bcM_{1,10}} &= x^{10} + 1014x^9 + 66090x^8 + 920263x^7 + 3975949x^6 + 6349238x^5 + 3975949x^4 + 920263x^3 + \cdots.
%66090x^2 + 1014x + 1
\end{align*}
\end{exam}
Note that this matches perfectly with \cite[eM1bar.txt]{BeT} and \cite[page 491]{Getzler2}. 

\medskip

By Theorems \ref{17} and \ref{22}, it is straightforward to write up a computer program that computes the Hodge-Deligne and Poincar\'e polynomials of $\bcM_{1,n}$ for higher $n$, quite effeciently. For $n=50$, we get the following picture for the probability distribution \beq\label{31}
p_{n,i}=\frac{\dim H^{2i}(\bcM_{1,n})}{\sum_{k=0}^n\dim H^{2k}(\bcM_{1,n})}.\eeq
The red curve in the picture is the graph of a Gaussian function. 

\begin{center}
\includegraphics[width=10cm]{image.png}
\end{center}

In the subsequent paper \cite{CKg1}, we will prove that indeed the probability distribution \eqref{31} is asymptotically Gaussian. 
Furthermore, we will prove a formula that gives us effective estimates of the Betti numbers for $n$ large and show that the Betti sequence is asymptotically 3-ultra log concave. 


\bigskip

\bibliographystyle{amsplain}
\begin{thebibliography}{99}
%	\bibitem{ACM} P. Aluffi, S. Chen and M. Marcolli. {\em Log concavity of the Grothendieck class of $\overline{\mathcal{M}}_{0,n}$}. J. Algebraic Combin. \textbf{61} (4) (2025), 47.
    \bibitem{BK} G. B\'erczi and Y.-H. Kiem. {\em Real-rootnedness of the Poincare polynomials of $\overline{\mathcal{M}}_{0,n}$ - an AI assisted proof.} Preprint, arXiv:2605.29151.
    \bibitem{BKs} G. B\'erczi and Y.-H. Kiem. {\em Stratified motivic invariants and bivariate deformations of Poincar\'e polynomials.} Preprint, arXiv:2607.23890.
    \bibitem{BeT} J. Bergstr\"om. Table of cohomology of $\bcM_{1,n}$.  https://github.com/jonasbergstroem/Cohomology-of-moduli-spaces-of-curves/blob/main/eM1nbar.txt
%    \bibitem{CK0} J. Choi and Y.-H. Kiem. {\em A new method toward the Landau-Ginzburg/Calabi-Yau correspondence via quasi-maps.} Math. Z. 294 (2020), 161–199.
    \bibitem{CK} J. Choi and Y.-H. Kiem. {\em Asymptotic distribution of the Betti numbers of $\bcM_{0,n}$}. Proc. Natl. Acad. Sci. USA (2026) 123, No. 18, e2601111123.
    \bibitem{CKg1} J. Choi and Y.-H. Kiem. {\em Asymptotic distribution of the Betti numbers of $\bcM_{1,n}$}. In preparation. 
    \bibitem{CKL1} J. Choi, Y.-H. Kiem and D. Lee. {\em Representations on the cohomology of $\overline{\mathcal{M}}_{0,n}$.} Adv. Math. \textbf{435} (2023), 109364.
    \bibitem{CKL2} J. Choi, Y.-H. Kiem and D. Lee. {\em Recursive algorithm and log-concavity of representations on the cohomology of $\bcM_{0,n}$}. Preprint, arXiv:2408.10728.
    \bibitem{CKL3} J. Choi, Y.-H. Kiem and D. Lee. {\em Characteristic polynomial of $\bcM_{0,n}$ and log-concavity}. Proc. Lond. Math. Soc. (3) 131 (2025), no. 5, Paper No. e70096.
    \bibitem{DM} P. Deligne and D. Mumford. {\em The irreducibility of the space of curves of given genus}. Publ. Math. Inst. Hautes Études Sci. (36) (1969) 75–109.
  %  \bibitem{FS} P. Flajolet and R. Sedgewick. {\em Analytic combinatorics}. Cambridge Univ. Press, Cambridge, 2009.
%    \bibitem{FM} W. Fulton and R. MacPherson. {\em A compactification of configuration spaces}. Ann. of Math. (2) 139 (1994), no. 1, 183-–225.
%    \bibitem{Getzler1} E. Getzler. {\em Operads and moduli spaces of genus zero Riemann surfaces.} The moduli space of curves, Progr. Math., vol. 129, Birkh\"auser, 1995, 199--230.
    \bibitem{Getzler2} E. Getzler. {\em The semi-classical approximation for modular operads.} Comm. Math. Phys. 194 (1998), 481--491.
    \bibitem{Getzler3} E. Getzler. {\em Resolving mixed Hodge modules on configuration spaces.} Duke Math. J. 96 (1999), 175--203.
    \bibitem{Hst} B. Hassett. {\em Moduli spaces of weighted pointed stable curves.} Adv. Math. 173 (2003) 316–352.
    \bibitem{KSY} S. Kannan, S. Serpente and C. H. Yun. {\em Equivariant Hodge polynomials of heavy/light moduli spaces.} Forum Math Sigma 12 (2024), e34, 1--19. 
%    \bibitem{Kapra} M. Kapranov. {\em Chow quotients of Grassmannians, I.} I.M. Gel’fand seminar, Adv. Soviet Math., vol. 16, Part 2, Amer. Math. Soc., Providence, RI, 1993, 29--110.
    \bibitem{Keel} S. Keel. {\em Intersection theory on the moduli space of stable n-pointed curves of genus zero}. Trans. Am. Math. Soc. \textbf{330} (2) (1992) 545--574.
%    \bibitem{KM} Y.-H. Kiem and H.-B. Moon. {\em Moduli space of weighted pointed stable rational curves via GIT}. Osaka J. Math. \textbf{48} (2011) 1115–-1140.
%    \bibitem{Kir} F. C. Kirwan. {\em Cohomology of quotients in symplectic and algebraic geometry.} Princeton University Press, 1984.
	\bibitem{Knu} F. Knudsen. {\em The projectivity of the moduli space of stable curves, II}. Math. Scand. \textbf{52} (1983), 161--199.
    \bibitem{Manin} Y. Manin. {\em Frobenius Manifolds, Quantum Cohomology, and Moduli Spaces}. AMS Colloquium Publ. 47, Providence, RI
    \bibitem{Shi} G. Shimura. {\em Introduction to the arithmetic theory of automorphic functions.} Princeton University Press, 1971.
\end{thebibliography}


\end{document}


\subsection{Matching with \cite{BeT}}
The non-equivariant Poincar\'e polynomial $P_{\overline{\mathcal{M}}_{1,n}}(L)$ can be independently computed from the $S_n$-equivariant Euler characteristics by substituting each Schur function $s_\lambda$ with the dimension $f^\lambda$ of its corresponding irreducible representation. This dimension is calculated using the hook-length formula:
\[
f^\lambda = \frac{n!}{\prod_{c \in \lambda} h_c}
\]
For example, for $n=3$, the partitions of 3 have dimensions $f^{[3]} = 1$ and $f^{[2,1]} = \frac{3!}{3 \cdot 1 \cdot 1} = 2$. Applying these to the equivariant data yields:
\[
P_{\overline{\mathcal{M}}_{1,3}}(L) = 1 \cdot (L^3+3L^2+3L+1) + 2 \cdot (L^2+L) = L^3 + 5L^2 + 5L + 1
\]
Evaluating $P_{\overline{\mathcal{M}}_{1,n}}(L) = \sum_{\lambda \vdash n} P_\lambda(L) f^\lambda$ systematically for all $n \le 10$ yields exactly the sequence derived from our recurrence:
\begin{align*}
P_{\overline{\mathcal{M}}_{1,1}}(L) &= L + 1 \\
P_{\overline{\mathcal{M}}_{1,2}}(L) &= L^2 + 2L + 1 \\
P_{\overline{\mathcal{M}}_{1,3}}(L) &= L^3 + 5L^2 + 5L + 1 \\
P_{\overline{\mathcal{M}}_{1,4}}(L) &= L^4 + 12L^3 + 23L^2 + 12L + 1 \\
P_{\overline{\mathcal{M}}_{1,5}}(L) &= L^5 + 27L^4 + 102L^3 + 102L^2 + 27L + 1 \\
P_{\overline{\mathcal{M}}_{1,6}}(L) &= L^6 + 58L^5 + 421L^4 + 756L^3 + 421L^2 + 58L + 1 \\
P_{\overline{\mathcal{M}}_{1,7}}(L) &= L^7 + 121L^6 + 1612L^5 + 5077L^4 + 5077L^3 + 1612L^2 + 121L + 1 \\
P_{\overline{\mathcal{M}}_{1,8}}(L) &= L^8 + 248L^7 + 5802L^6 + 31072L^5 + 52402L^4 + 31072L^3 + 5802L^2 + 248L + 1 \\
P_{\overline{\mathcal{M}}_{1,9}}(L) &= L^9 + 503L^8 + 19925L^7 + 175036L^6 + 480097L^5 + 480097L^4 + 175036L^3 + 19925L^2 + 503L + 1 \\
P_{\overline{\mathcal{M}}_{1,10}}(L) &= L^{10} + 1014L^9 + 66090L^8 + 920263L^7 + 3975949L^6 + 6349238L^5 + 3975949L^4 + 920263L^3 + \cdots. % 66090L^2 + 1014L + 1
\end{align*}

Consequently, the analytic difference between the generated sequence and the geometric invariants rigorously evaluates to zero:
\[
P_{\overline{\mathcal{M}}_{1,n}}(x) - P_{n,0}(x) = 0 \quad \text{for all } 1 \le n \le 10.
\]
